\begin{table}[htbp!]
\centering
\caption{Ambiente de execução de cada rodada --- referência em nuvem (Groq)}
\label{tab:ambiente_groq}
\footnotesize
\ajustartabela{%
\begin{tabular}{|l|l|l|c|c|c|c|}
\hline
\textbf{Modelo} & \textbf{GPU / máquina} & \textbf{Servidor} & \textbf{Rep.} & \textbf{Exec.} & \textbf{Na GPU} & \textbf{VRAM livre} \\
\hline
GPT-OSS 20B (Groq) & servidores do provedor & Groq (API) & 1 & 310 & --- & --- \\
GPT-OSS 120B (Groq) & servidores do provedor & Groq (API) & 1 & 310 & --- & --- \\
\hline
\end{tabular}}
\fonte{Elaborado pelo autor a partir de \texttt{results/<modelo>/manifesto.json}. Na GPU: parcela do modelo carregado (pesos e contexto) alocada na VRAM pelo Ollama, medida após o aquecimento; VRAM livre: no início da rodada, com os demais modelos descarregados. Data de referência: 2026-09-24. Python 3.14.4; langchain-core 1.6.4, langchain-ollama 1.1.0, langchain-openai 1.6.6.}
\end{table}
